%\documentclass[notes,usenames,dvipsnames]{beamer}
%\documentclass[notes]{beamer}       % print frame + notes
%\documentclass[notes=only]{beamer}   % only notes
\documentclass[usenames,dvipsnames]{beamer}             % only frames
\usepackage{pgfpages}
%\setbeameroption{show notes on second screen}
%\setbeameroption{show notes}
%subfigures
\usepackage{caption}
\usepackage{subcaption}

%tables packages
\usepackage{multirow}

% math
\usepackage{amsmath}

% bash command
\usepackage{graphicx}
\usepackage{listings}
% varbatim for ascii figures
\usepackage[T1]{fontenc}
\usepackage[utf8]{inputenc}
%\usepackage{verbatim}
\usepackage{lipsum} % for context
\usepackage{fancyvrb}
\usepackage{varwidth}


\newsavebox{\asciigcn}

% notes prefixed in pympress
\addtobeamertemplate{note page}{}{\thispdfpagelabel{notes:\insertframenumber}}
%theme used
\usetheme{Madrid}
%Information to be included in the title page:
\title[Computer Architecture] %optional
{Arhitectura Calculatoarelor}

\subtitle{Seria AB - Curs 1 - Structură Calculator Numeric}

\author[Ștefan-Dan Ciocîrlan] % (optional, for multiple authors)
{}

\institute[NUSTPB] % (optional)
{
  \inst{}%
  National University of Science and Technology\\
  POLITEHNICA Bucharest
}

\date[NUSTPB 2025] % (optional)
{Computer Architecture}

\logo{
\includegraphics[height=1.3cm]{../../media/logo-poli-color9.png}
%\includegraphics[height=0.7cm]{media/sbipLogo.pdf}
}


% Roman numerals
\newcommand*{\rom}[1]{\expandafter\@slowromancap\romannumeral #1@}
%split
\usepackage{amsmath}
%colors
%\usepackage[usenames,dvipsnames]{color} %loaded by the dcoument class
%subfigure
\usepackage{subcaption}
% block over block uncover
\setbeamercovered{invisible}
%\setbeamercovered{transparent}

%extra slide content
\AtBeginSection[]
{
  \begin{frame}
    \frametitle{Conținut}
    \tableofcontents[currentsection]
  \end{frame}
}
%notes or not
%\setbeamertemplate{note page}[plain]

\begin{document}

\frame{\titlepage}


\section{Calculator - Sistem de prelucrare a informației}
\begin{frame}
    \frametitle{ALU Structure}
    \begin{columns}
        \column{0.5\textwidth}
        \begin{itemize}
            \item ALU (Arithmetic Logic Unit) is a combinational circuit.
            \item It performs arithmetic and logic operations.
            \item There is a control input for selecting the desired operation.
            \item It might include status flags.
            \item The ALU can vary in cycle time and complexity.
        \end{itemize}

        \column{0.5\textwidth}
            \newsavebox{\asciialu}
            \begin{lrbox}{\asciialu}
                \begin{varwidth}{\maxdimen}
                \VerbatimInput[fontsize=\scriptsize]{media/alu.ascii}
                \end{varwidth}
            \end{lrbox}%

            \begin{figure}[h]
                \centering
                \scalebox{0.7}{\usebox{\asciialu}}
            \end{figure}
    \end{columns}
    \note{
    }
\end{frame}


\section{Model teoretic al Calculatorului Numeric}
\begin{frame}
    \frametitle{Expersii Regulate/Logică combinațională}
    \newsavebox{\asciiregular}
    \begin{lrbox}{\asciiregular}
        \begin{varwidth}{\maxdimen}
        \VerbatimInput[fontsize=\scriptsize]{media/regular.ascii}
        \end{varwidth}
    \end{lrbox}%

    \begin{figure}[h]
        \centering
        \scalebox{0.8}{\usebox{\asciiregular}}
    \end{figure}
    \begin{itemize}
        \item Expresii regulate $O = [a-zA-Z0-9]+$
        \item Logică combinațională $O = i_0 \oplus i_1$
        \item Are nevoie de toată intrarea pentru a produce ieșirea
    \end{itemize}
\end{frame}

\begin{frame}
    \frametitle{Automat Finite}
    \newsavebox{\asciifsm}
    \begin{lrbox}{\asciifsm}
        \begin{varwidth}{\maxdimen}
        \VerbatimInput[fontsize=\scriptsize]{media/fsm.ascii}
        \end{varwidth}
    \end{lrbox}%

    \begin{figure}[h]
        \centering
        \scalebox{0.8}{\usebox{\asciifsm}}
    \end{figure}
    \begin{itemize}
        \item Are o stare internă ce poate fi modificată
        \item Nu poate accessa intrări anterior anterior
    \end{itemize}
\end{frame}


\begin{frame}
    \frametitle{Automat cu stivă/pushdown}
    \newsavebox{\asciistackfsm}
    \begin{lrbox}{\asciistackfsm}
        \begin{varwidth}{\maxdimen}
        \VerbatimInput[fontsize=\scriptsize]{media/stackfsm.ascii}
        \end{varwidth}
    \end{lrbox}%

    \begin{figure}[h]
        \centering
        \scalebox{0.8}{\usebox{\asciistackfsm}}
    \end{figure}
    \begin{itemize}
        \item Are o stivă unde poate stoca și re accesa intrările.
        \item Stiva este limitată, o intrare citită salvată și folosită nu poate fi folosită din nou.
    \end{itemize}
\end{frame}


\begin{frame}
    \frametitle{Model Turing Machine}
    \newsavebox{\asciituring}
    \begin{lrbox}{\asciituring}
        \begin{varwidth}{\maxdimen}
        \VerbatimInput[fontsize=\scriptsize]{media/turing.ascii}
        \end{varwidth}
    \end{lrbox}%

    \begin{figure}[h]
        \centering
        \scalebox{0.8}{\usebox{\asciituring}}
    \end{figure}
    \begin{itemize}
        \item Are o rolă infinită și poate reaccesa intrările și ieșirle anterioare.
    \end{itemize}
\end{frame}

\section{Model structural al Calculatorului Numeric}
\begin{frame}
    \frametitle{Proprietați CN}
\begin{itemize}
    \item Determinist în teorie, predictibil în inginerie cu eroare neglijabilă ce poate fi modelată.
    \item Modelat după Turing Machine.
\end{itemize}

\end{frame}

\begin{frame}
    \frametitle{Model general de CN}
    \begin{lrbox}{\asciigcn}
        \begin{varwidth}{\maxdimen}
        \VerbatimInput[fontsize=\scriptsize]{media/gcn.ascii}
        \end{varwidth}
    \end{lrbox}%

    \begin{figure}[h]
        \centering
        \scalebox{0.9}{\usebox{\asciigcn}}
    \end{figure}

    \begin{itemize}
        \item Diferențe între reprezentarea internă și externă a datelor.
    \end{itemize}
\end{frame}


\begin{frame}
    \frametitle{Model structural al CN}
    \newsavebox{\asciimscn}
    \begin{lrbox}{\asciimscn}
        \begin{varwidth}{\maxdimen}
        \VerbatimInput[fontsize=\scriptsize]{media/mscn.ascii}
        \end{varwidth}
    \end{lrbox}%

    \begin{figure}[h]
        \centering
        \scalebox{0.8}{\usebox{\asciimscn}}
    \end{figure}
    \begin{itemize}
        \item Diferneță între echipamentele periferice și calculatorul numeric.
        \item Memorie internă și unitatea centrală de procesare.
    \end{itemize}
\end{frame}


\begin{frame}
    \frametitle{Structura unui CN}
    \newsavebox{\asciistructcn}
    \begin{lrbox}{\asciistructcn}
        \begin{varwidth}{\maxdimen}
        \VerbatimInput[fontsize=\scriptsize]{media/structcn.ascii}
        \end{varwidth}
    \end{lrbox}%

    \begin{figure}[h]
        \centering
        \scalebox{0.65}{\usebox{\asciistructcn}}
    \end{figure}
    \begin{itemize}
        \item Ierarhie de memorie.
        \item Control și procesare.
    \end{itemize}
\end{frame}


\begin{frame}
    \frametitle{Structura unui Sistem de Calcul (SC)}
    \newsavebox{\asciistructsoc}
    \begin{lrbox}{\asciistructsoc}
        \begin{varwidth}{\maxdimen}
        \VerbatimInput[fontsize=\scriptsize]{media/structsoc.ascii}
        \end{varwidth}
    \end{lrbox}%

    \begin{figure}[h]
        \centering
        \scalebox{0.47}{\usebox{\asciistructsoc}}
    \end{figure}
    \begin{itemize}
        \item Control memorie.
    \end{itemize}
\end{frame}

\section{Model funcțional al Calculatorului Numeric}
\begin{frame}
    \frametitle{Niveluri funcționale ale unui CN}
\begin{enumerate}
    \item Dispozitve și circuite electronice (Transiztoare, portți logice) - hardware
    \item Unități funcționale (UAL, Memorie, Interfețe) - firmware (Micropgramare) + hardware
    \item Mașină fizică Strcuturi de interconectare și echipamente periferice - hardware
    \item Nucleul sistem de operare - BIOS - firmware
    \item Sistem de Operare și API-ul său - software
    \item Limbaj de programare și compilator - software
    \item Aplicații și biblioteci de programare - software
    \item Interfața cu utilizatorul - software
\end{enumerate}

\end{frame}

\begin{frame}
    \frametitle{Hardware vs Software}
\begin{itemize}
    \item Anumite funcții pot fi implementate atât în hardware cât și în software
    \item Decizia de a implementa o funcție în hardware sau software se bazează pe:
    \begin{itemize}
        \item Costul implementării
        \item Viteza de execuție
        \item Raport cost/performanță
        \item Siguranță în funcționare
        \item Frecvența unor modificări
    \end{itemize}
    \item Tendința actuală este de a trece cât mai multe funcții în hardware
\end{itemize}

\end{frame}

\section{Unități funcționale a Calculatorului Numeric}
\begin{frame}
    \frametitle{Memoria}
    Tipuri de memorie:
    \begin{itemize}
        \item Regiștri
        \item Tampon (Cache)
        \item Principală (RAM)
        \item Memorie auxiliară/externă (HDD, SSD) - Poate să fie parte din subsistemul de intrare/ieșire
    \end{itemize}
    Proprietați:
    \begin{itemize}
        \item Fiecare Locație are același număr de biți
        \item Fiecare locație are o adresă unică 
        \item Spațiul de adresare este omoegen și toate locațiile de memorie sunt echivalente
    \end{itemize}
\end{frame}

\begin{frame}
    \frametitle{Memoria}
    \newsavebox{\asciimemwrite}
    \begin{lrbox}{\asciimemwrite}
        \begin{varwidth}{\maxdimen}
        \VerbatimInput[fontsize=\scriptsize]{media/memwrite.ascii}
        \end{varwidth}
    \end{lrbox}%

    \begin{figure}[h]
        \centering
        \scalebox{0.8}{\usebox{\asciimemwrite}}
    \end{figure}
    \begin{itemize}
        \item Timp acces memorie depinde de lungimea addresei
        \item Timp ciclu de scriere/citire depinde de lungimea cuvântului
    \end{itemize}
\end{frame}


\begin{frame}
    \frametitle{Unitate centrală de prelucrare}
    Componente:
    \begin{itemize}
        \item Unitatea aritmetică și logică (UAL).
        \item Unitatea de comandă.
    \end{itemize}
    Tipuri de UAL:
    \begin{itemize}
        \item UAL cu o singură magistrală.
        \item UAL cu o singură magistrală și un acumulator.
        \item UAL cu trei magistrale.
    \end{itemize}
\end{frame}

\begin{frame}
    \frametitle{Unitatea de comandă}
    \begin{itemize}
        \item Decodifică instrucțiunile și generează semnalele de control.
        \item Preia instrucțiunea de la adresa dată de registru contor program (CP).
        \item Poate să fie implementată:
        \begin{itemize}
            \item Unitate de comandă microprogramată
            \item Unitate de comandă convențională (Automat)
        \end{itemize}
        \item Dependent de setul de instrucțiuni.
    \end{itemize}
\end{frame}

\begin{frame}
    \frametitle{Set de instrucțiuni}
    Tipuri de instrucțiuni:
    \begin{itemize}
        \item Transfer de date
        \item Operații aritmetice
        \item Operații logice și de deplasare
        \item Operații de comparare
        \item Control de flux
    \end{itemize}
    Fiecare instrucțiune are un cod de operație și specificatori de operanzi.
\end{frame}

\begin{frame}
    \frametitle{Subsistem de intrare/ieșire}
    Rolul subsistemului de intrare/ieșire este de a asigura conversia de format și viteză între UCP și dispozitivele periferice.
\end{frame}

\begin{frame}
    \frametitle{Limbaje de programare}
    \newsavebox{\asciiplcomp}
    \begin{lrbox}{\asciiplcomp}
        \begin{varwidth}{\maxdimen}
        \VerbatimInput[fontsize=\scriptsize]{media/plcomp.ascii}
        \end{varwidth}
    \end{lrbox}%

    \begin{figure}[h]
        \centering
        \scalebox{0.8}{\usebox{\asciiplcomp}}
    \end{figure}
    Rolul limbajelor de programare este de a crea iluzia unei ierarhi de calculatoare virtuale.
\end{frame}

\section{Q\&A}
\begin{frame}
\end{frame}

%\begin{frame}
%\frametitle{Table of Contents}
%\tableofcontents
%\end{frame}


\end{document}